\documentclass[11pt]{article}

% ============================================================
% Packages
% ============================================================
\usepackage[margin=1in]{geometry}
\usepackage{amsmath,amssymb,mathtools,amsthm}
\usepackage{enumitem}
\usepackage{hyperref}
\usepackage{xcolor}
\usepackage{microtype}
\usepackage{tikz}

% ============================================================
% Macros
% ============================================================
\newcommand{\R}{\mathbb{R}}
\newcommand{\Pbb}{\mathbb{P}}
\newcommand{\Qbb}{\mathbb{Q}}
\newcommand{\E}{\mathbb{E}}
\newcommand{\dd}{\,\mathrm{d}}

% ============================================================
% Environments
% ============================================================
\theoremstyle{plain}
\newtheorem{theorem}{Theorem}[section]
\newtheorem{proposition}[theorem]{Proposition}

\theoremstyle{definition}
\newtheorem{definition}[theorem]{Definition}
\newtheorem{example}[theorem]{Example}

\theoremstyle{remark}
\newtheorem{remark}[theorem]{Remark}

% ============================================================
% Title
% ============================================================
\title{Option Pricing --- Lecture Notes}
\author{Nathan Sauldubois}
\date{MSFE6233 --- Local Volatility, Stochastic Volatility, and Stochastic Interest Rates}

\begin{document}

\maketitle
\vspace{-1em}

\noindent\textbf{Lecture title.}
\emph{Beyond Black--Scholes: local volatility, stochastic volatility, and stochastic interest rates}

\vspace{0.5em}
\hrule
\vspace{1em}

\noindent\textbf{Goals of the lecture.}
\begin{itemize}[leftmargin=1.4em]
    \item understand why the Black--Scholes model is too rigid to fit market smiles and skews;
    \item introduce local volatility models and their link with the implied volatility surface;
    \item present the main ideas behind stochastic volatility models;
    \item introduce stochastic interest-rate modeling and its implications for pricing fixed-income securities.
\end{itemize}

\tableofcontents
\bigskip

% ============================================================
\section{Introduction and motivation}
% ============================================================

\subsection{Limitations of the Black--Scholes model}

The Black--Scholes model is one of the cornerstones of modern option
pricing.
Its success comes from its mathematical tractability, the existence of
closed-form pricing formulas for European options, and the fact that it
provides a very clear benchmark for hedging and risk management.  

However, the model relies on two strong assumptions:
\begin{itemize}[leftmargin=1.4em]
    \item the volatility of the underlying asset is constant,
    \item the interest rate is deterministic, and often even constant.
\end{itemize}

These assumptions are convenient mathematically, but they are not
consistent with real market data.
In particular, if the stock price satisfies
\[
\frac{\dd S_t}{S_t}=\mu\,\dd t+\sigma\,\dd W_t,
\]
with constant \(\sigma\), then log-returns are Gaussian and their
distribution is entirely characterized by a fixed variance parameter.  
Empirically, this is not what one observes in financial markets.
Historical returns exhibit important deviations from the Gaussian model:
\begin{itemize}[leftmargin=1.4em]
    \item volatility varies through time,
    \item extreme events occur more often than predicted by a Gaussian law,
    \item distributions of returns are often skewed and leptokurtic.
\end{itemize}

In particular, financial data are known to display \emph{fat tails},
meaning that extreme movements are more likely than under the normal
distribution.
This is often measured through the kurtosis, which is strictly larger
than the Gaussian benchmark value \(3\). 
Thus, even though the Black--Scholes model remains extremely useful in
practice, it is now well understood that it cannot be a fully realistic
description of financial markets. 

\subsection{Volatility smile, skew, and term structure}

One of the clearest empirical failures of the Black--Scholes model comes
from the behavior of implied volatility.

Recall that, in the Black--Scholes framework, one can invert the pricing
formula and define the implied volatility \(\sigma_{\mathrm{imp}}(K,T)\)
as the volatility parameter that reproduces the observed market option
price.
If the Black--Scholes model were correct, implied volatility would be
constant across strikes and maturities.  

In real markets, this is not the case.
Instead, implied volatility depends on both the strike and the maturity,
leading to a full \emph{implied volatility surface}.  

Two important patterns are observed:
\begin{itemize}[leftmargin=1.4em]
    \item the \emph{volatility smile}, where implied volatility is higher
    for deep in-the-money and deep out-of-the-money options than near the
    money,
    \item the \emph{volatility skew}, especially visible in equity markets,
    where implied volatility tends to increase as the strike decreases.  
\end{itemize}

The volatility skew reflects the fact that market participants attach
significant value to downside protection.
Out-of-the-money put options are in strong demand because they insure
against large downward moves, and this increased demand pushes their
prices upward.
When translated into Black--Scholes implied volatilities, this appears
as a higher implied volatility for low strikes. 

In addition, implied volatility also varies with maturity.
This \emph{term structure of volatility} shows that the market does not
price short-dated and long-dated uncertainty in the same way.
Hence a single constant volatility parameter cannot reproduce the whole
cross-section of market option prices. 

\subsection{Why go beyond constant volatility and constant rates?}

These empirical observations naturally motivate richer models.

A first direction is to replace the constant volatility \(\sigma\) by a
deterministic function of time and spot,
\[
\sigma_{\mathrm{loc}}(t,S_t),
\]
leading to a \emph{local volatility model}.
Such models are designed precisely to fit the observed implied volatility
surface at a given date. 

A second direction is to model volatility itself as a random process.
This leads to \emph{stochastic volatility models}, which are able to
capture richer smile dynamics and more realistic time-series behavior of
returns.
They are also closely related to the observed clustering of volatility
in financial data. 

Finally, the assumption of deterministic interest rates is often too
simplistic, especially when pricing fixed-income products or derivatives
with significant exposure to the term structure of rates.
This motivates the study of \emph{stochastic interest-rate models},
where bond prices, yields, and forward rates become the main state
variables.

From a practical point of view, all these extensions are driven by the
same need:
\begin{itemize}[leftmargin=1.4em]
    \item reproduce market prices more accurately,
    \item obtain models consistent with observed implied volatility surfaces,
    \item improve hedging and risk management,
    \item and build frameworks suitable for calibration to liquid market data.  
\end{itemize}

\subsection{Roadmap of the lecture}

This lecture is organized around three main modeling extensions beyond
Black--Scholes.

We begin with \emph{local volatility models}.
The starting point is the implied volatility surface observed in the
market.
We introduce local volatility as a deterministic function
\(\sigma_{\mathrm{loc}}(t,s)\), derive the corresponding pricing PDE,
and present Dupire's formula, which links local volatility to European
call prices observed across strikes and maturities.  

We then turn to \emph{stochastic volatility models}.
The goal is to go one step further and allow volatility itself to be
random.
We present the general structure of such models, explain their economic
motivation, and discuss the role of additional factors such as the
correlation between the asset and its volatility.

Finally, we introduce \emph{stochastic interest-rate models}.
We recall the basic objects of fixed-income modeling, such as
zero-coupon bonds, yields, and forward rates, and explain why the
evolution of the whole term structure must often be modeled explicitly.
This leads to factor models and, at a more advanced level, to the


The general philosophy of the lecture is the following:
\begin{itemize}[leftmargin=1.4em]
    \item identify the empirical shortcomings of the Black--Scholes model,
    \item introduce richer but still tractable alternatives,
    \item and understand how these models connect market observables,
    such as option prices or yield curves, to the dynamics of the
    underlying sources of risk.
\end{itemize}

% ============================================================
\section{Local volatility models}
% ============================================================

\subsection{Definition of the local volatility model}

A first way to go beyond the Black--Scholes model is to keep the market
complete and Markovian, but replace the constant volatility parameter by
a deterministic function of time and spot. This leads to the
\emph{local volatility model}. 

\paragraph{Definition.}
Under the risk-neutral measure, the asset price process \(S\) is assumed
to satisfy
\[
\mathrm{d}S_t
=
rS_t\,\mathrm{d}t
+
\sigma_{\mathrm{loc}}(t,S_t)\,S_t\,\mathrm{d}W_t,
\]
or, equivalently,
\[
\frac{\mathrm{d}S_t}{S_t}
=
r\,\mathrm{d}t
+
\sigma_{\mathrm{loc}}(t,S_t)\,\mathrm{d}W_t.
\]
Here, the function
\[
\sigma_{\mathrm{loc}}:[0,T]\times \R_+ \to \R_+
\]
is called the \emph{local volatility function}. It is a deterministic
function of the current time and the current asset price. 

\paragraph{Interpretation.}
The terminology ``local'' reflects the fact that the instantaneous
volatility at time \(t\) depends only on the current state \(S_t\), and
not on the past trajectory of the process or on any additional hidden
factor. Thus the model remains one-factor and Markovian. 

\paragraph{Comparison with Black--Scholes.}
The Black--Scholes model is recovered as the special case
\[
\sigma_{\mathrm{loc}}(t,s)\equiv \sigma,
\]
where the volatility is constant.
The local volatility model is therefore a natural extension of
Black--Scholes obtained by allowing the volatility coefficient to vary
with both time and spot. 

\paragraph{Why this model is introduced.}
The main motivation is calibration.
A single constant volatility parameter cannot reproduce the full market
surface of European option prices across strikes and maturities.
By allowing \(\sigma_{\mathrm{loc}}(t,s)\) to vary with \(t\) and \(s\),
one obtains a model with enough flexibility to fit the observed implied
volatility surface at a given date. 

\paragraph{Pricing under local volatility.}
Since the dynamics remain Markovian under the risk-neutral measure, the
price \(C(t,s)\) of a European claim with payoff \(g(S_T)\) is still
given by a risk-neutral expectation
\[
C(t,s)
=
\E^\Qbb\!\left[
e^{-r(T-t)}g(S_T)
\;\middle|\;
S_t=s
\right],
\]
and it solves the valuation PDE
\[
\partial_t C
+
rs\,\partial_s C
+
\frac12 \sigma_{\mathrm{loc}}(t,s)^2 s^2\,\partial_{ss}C
-
rC
=
0,
\qquad
C(T,s)=g(s).
\]
Thus the local volatility model preserves much of the analytic structure
of Black--Scholes while allowing much richer implied-volatility
patterns. 

\subsection{Risk-neutral dynamics under local volatility}

We now describe more precisely the dynamics of the asset price under a
local volatility model when we work under the risk-neutral measure.

\paragraph{Risk-neutral specification.}
As in the Black--Scholes model, we assume that under the risk-neutral
measure \(\Qbb\), the drift of the asset price is the short rate.
The difference is that the volatility is no longer constant, but depends
on both time and the current level of the asset.

Thus the dynamics are
\[
\mathrm{d}S_t
=
rS_t\,\mathrm{d}t
+
\sigma_{\mathrm{loc}}(t,S_t)\,S_t\,\mathrm{d}W_t,
\]
or equivalently
\[
\frac{\mathrm{d}S_t}{S_t}
=
r\,\mathrm{d}t
+
\sigma_{\mathrm{loc}}(t,S_t)\,\mathrm{d}W_t.
\]

\paragraph{Comparison with Black--Scholes.}
In the Black--Scholes model, the volatility coefficient is constant, so
the diffusion term is simply
\[
\sigma S_t\,\mathrm{d}W_t.
\]
In the local volatility model, this is replaced by
\[
\sigma_{\mathrm{loc}}(t,S_t)S_t\,\mathrm{d}W_t.
\]
Hence the instantaneous variance of the asset return is now
\[
\sigma_{\mathrm{loc}}(t,S_t)^2.
\]

\paragraph{Markovian structure.}
A key feature of the model is that the dynamics remain one-dimensional
and Markovian.
Indeed, the future evolution of the process depends only on the current
time \(t\) and the current value \(S_t\).
No additional factor is introduced.

This is one of the main differences with stochastic volatility models,
where volatility itself is an extra random process.

\paragraph{Log-price dynamics.}
Applying It\^o's formula to \(\log S_t\), we obtain
\[
\mathrm{d}\log S_t
=
\Big(
r-\frac12\sigma_{\mathrm{loc}}(t,S_t)^2
\Big)\mathrm{d}t
+
\sigma_{\mathrm{loc}}(t,S_t)\,\mathrm{d}W_t.
\]

Thus, unlike in the Black--Scholes model, the log-price is no longer
Gaussian in general, since its diffusion coefficient depends on the
current state of the process.

\paragraph{Discounted dynamics.}
If we define the discounted stock price by
\[
\widetilde S_t:=e^{-rt}S_t,
\]
then It\^o's product rule gives
\[
\mathrm{d}\widetilde S_t
=
e^{-rt}\sigma_{\mathrm{loc}}(t,S_t)S_t\,\mathrm{d}W_t.
\]
Hence
\[
\widetilde S_t
\]
is a local martingale under \(\Qbb\), and under suitable integrability
assumptions it is a true martingale.
This is the risk-neutral counterpart of the no-arbitrage condition.

\paragraph{Instantaneous covariance structure.}
Since the model is still driven by a single Brownian motion, the only
source of randomness is the asset itself.
The local volatility function changes the amplitude of the noise, but it
does not introduce any additional independent factor.

\paragraph{Economic interpretation.}
The local volatility function
\[
\sigma_{\mathrm{loc}}(t,s)
\]
can be interpreted as the market-implied instantaneous volatility of the
asset when the time is \(t\) and the spot level is \(s\).
In particular:
\begin{itemize}[leftmargin=1.4em]
    \item dependence on \(t\) captures the term structure of volatility,
    \item dependence on \(s\) captures the smile or skew effect across strikes.
\end{itemize}

\paragraph{Main lesson.}
The risk-neutral dynamics under local volatility keep the same general
structure as Black--Scholes,
\[
\text{drift }=rS_t,
\qquad
\text{diffusion }=\sigma_{\mathrm{loc}}(t,S_t)S_t,
\]
but the state-dependent volatility coefficient allows the model to fit a
much richer cross-section of option prices.

\subsection{Pricing PDE under a local volatility model}

We now derive the valuation equation satisfied by a European option in a
local volatility model.

\paragraph{Set-up.}
Assume that under the risk-neutral measure \(\Qbb\), the stock price
satisfies
\[
\mathrm{d}S_t
=
rS_t\,\mathrm{d}t
+
\sigma_{\mathrm{loc}}(t,S_t)S_t\,\mathrm{d}W_t,
\]
where \(r\) is constant and \(\sigma_{\mathrm{loc}}(t,s)\) is a
deterministic function of time and spot. 

Let \(g:\R_+\to\R\) be the payoff function of a European claim maturing
at time \(T\), and define its time-\(t\) price by
\[
C(t,s)
=
\E^\Qbb\!\left[
e^{-r(T-t)}g(S_T)
\;\middle|\;
S_t=s
\right].
\]
As in the Black--Scholes framework, the option price is characterized by
a backward parabolic PDE. 

\paragraph{Valuation equation.}
The pricing PDE takes the form
\[
\boxed{
\partial_t C(t,s)
+
rs\,\partial_s C(t,s)
+
\frac12 \sigma_{\mathrm{loc}}(t,s)^2 s^2\,\partial_{ss}C(t,s)
-
rC(t,s)
=
0,
}
\]
for \(0\le t<T\), together with the terminal condition
\[
\boxed{
C(T,s)=g(s).
}
\]
This is exactly the Black--Scholes PDE with the constant volatility
\(\sigma\) replaced by the local volatility function
\(\sigma_{\mathrm{loc}}(t,s)\). 

\paragraph{Derivation.}
There are two standard ways to obtain this equation.

The first one is to apply the Feynman--Kac formula to the risk-neutral
expectation above.
Since \(S\) is a Markov diffusion with generator
\[
\mathcal L_t \varphi(s)
=
rs\,\varphi'(s)
+
\frac12 \sigma_{\mathrm{loc}}(t,s)^2 s^2\,\varphi''(s),
\]
the discounted price function must satisfy
\[
\partial_t C+\mathcal L_t C-rC=0.
\]

The second one is the usual replication argument:
apply It\^o's formula to \(C(t,S_t)\), choose the delta hedge
\[
\Delta_t=\partial_s C(t,S_t),
\]
and impose that the hedged portfolio earns the risk-free rate.
This gives the same PDE. The local volatility case is therefore a direct
extension of the standard Black--Scholes hedging argument. 

\paragraph{Interpretation.}
The valuation equation shows that, under local volatility, the option
price still depends only on time and spot,
\[
C_t=C(t,S_t),
\]
so the model remains one-factor and Markovian.
In particular, European option pricing remains analytically tractable at
the PDE level, even though the model can fit a much richer cross-section
of market prices than Black--Scholes. 

\paragraph{Why this is not enough for calibration.}
A first naive idea would be to recover the local volatility function by
solving the pricing PDE for \(\sigma_{\mathrm{loc}}\), namely
\[
\sigma_{\mathrm{loc}}(t,s)^2
=
\frac{
rC(t,s)-\partial_t C(t,s)-rs\,\partial_s C(t,s)
}{
\frac12 s^2\,\partial_{ss}C(t,s)
}.
\]
However, this is not directly implementable in practice, because the
derivatives with respect to the spot variable \(s\) are not observed on
the market: market quotes are given across strikes and maturities, not
as a function of the current spot state variable. This is precisely why
Dupire's forward formulation in \((T,K)\) variables is needed.  

\subsection{Forward Kolmogorov equation and transition densities}

To connect local volatility models with market option prices, it is
useful to complement the backward pricing PDE with the corresponding
\emph{forward} equation satisfied by the law of the stock price.

\paragraph{Transition density.}
Assume that, under the risk-neutral measure, the stock price satisfies
\[
\mathrm{d}S_t
=
rS_t\,\mathrm{d}t
+
\sigma_{\mathrm{loc}}(t,S_t)S_t\,\mathrm{d}W_t.
\]
For \(0\le t<T\), we denote by
\[
p(t,s;T,y)
\]
the transition density of \(S_T\) conditional on \(S_t=s\), so that
\[
\Pbb(S_T\in \mathrm{d}y\mid S_t=s)=p(t,s;T,y)\,\mathrm{d}y.
\]

Equivalently, for any bounded measurable function \(g\),
\[
\E^\Qbb[g(S_T)\mid S_t=s]
=
\int_0^\infty g(y)\,p(t,s;T,y)\,\mathrm{d}y.
\]

\paragraph{Pricing formula in terms of the density.}
If \(C(t,s)\) denotes the price at time \(t\) of a European claim with
payoff \(g(S_T)\), then
\[
C(t,s)
=
e^{-r(T-t)}
\int_0^\infty g(y)\,p(t,s;T,y)\,\mathrm{d}y.
\]

In particular, for a European call option with strike \(K\),
\[
C(t,s;T,K)
=
e^{-r(T-t)}
\int_0^\infty (y-K)^+\,p(t,s;T,y)\,\mathrm{d}y.
\]

Thus the full cross-section of European option prices is determined by
the transition density of the stock under the risk-neutral measure.

\paragraph{Forward Kolmogorov equation.}
The transition density satisfies the forward Kolmogorov equation, also
called the \emph{Fokker--Planck equation}.
For fixed initial condition \((t,s)\), the function
\[
(T,y)\longmapsto p(t,s;T,y)
\]
solves
\[
\boxed{
\partial_T p(t,s;T,y)
=
-\partial_y\big(ry\,p(t,s;T,y)\big)
+
\frac12\,\partial_{yy}\!\Big(\sigma_{\mathrm{loc}}(T,y)^2y^2\,p(t,s;T,y)\Big).
}
\]

The initial condition is
\[
p(t,s;t,y)=\delta_s(y),
\]
in the sense of distributions.

\paragraph{Interpretation.}
This equation describes the time evolution of the probability density of
the stock price.
The first term corresponds to the transport generated by the drift
\(ry\), while the second term describes the diffusion induced by the
local volatility coefficient.

\paragraph{Backward PDE vs forward PDE.}
It is useful to distinguish clearly the two PDE viewpoints:
\begin{itemize}[leftmargin=1.4em]
    \item the \emph{backward} pricing PDE is satisfied by the option
    price \(C(t,s)\) as a function of current time and current spot;
    \item the \emph{forward} Kolmogorov equation is satisfied by the
    transition density \(p(t,s;T,y)\) as a function of maturity and
    future state variable.
\end{itemize}

Both equations encode the same diffusion model, but from two different
perspectives.

\paragraph{Why the forward equation matters.}
The forward formulation is particularly important for calibration,
because market option prices are quoted as functions of maturity \(T\)
and strike \(K\), not as functions of the current spot variable \(s\).
This makes the forward density approach better adapted to extracting the
local volatility surface from market data.

Indeed, Dupire's formula is obtained precisely by combining:
\begin{itemize}[leftmargin=1.4em]
    \item the representation of call prices in terms of the transition
    density,
    \item and the forward Kolmogorov equation satisfied by that density.
\end{itemize}

\paragraph{Key message.}
The local volatility model can be studied either through:
\begin{itemize}[leftmargin=1.4em]
    \item the backward pricing PDE for option prices,
    \item or the forward Kolmogorov equation for transition densities.
\end{itemize}
The forward point of view is the natural one when one wants to link the
model directly to the observed surface of European option prices.

\subsection{Dupire's formula}

\subsubsection{Statement of Dupire's formula}

We now state one of the central results of the local volatility
approach.

Suppose that, for each maturity \(T\) and strike \(K\), the market
provides the price
\[
C(T,K)
\]
of a European call option written on the underlying asset.
Assume that these call prices are sufficiently smooth in \((T,K)\), and
that the underlying asset follows a local volatility model under the
risk-neutral measure:
\[
\mathrm{d}S_t
=
rS_t\,\mathrm{d}t
+
\sigma_{\mathrm{loc}}(t,S_t)S_t\,\mathrm{d}W_t.
\]

Then the local volatility function is given by Dupire's formula
\[
\boxed{
\sigma_{\mathrm{loc}}(T,K)^2
=
\frac{
\partial_T C(T,K)+rK\,\partial_K C(T,K)
}{
\frac12 K^2\,\partial_{KK}C(T,K)
}.
}
\]

Equivalently,
\[
\boxed{
\partial_T C(T,K)
=
\frac12 \sigma_{\mathrm{loc}}(T,K)^2 K^2\,\partial_{KK}C(T,K)
-rK\,\partial_K C(T,K).
}
\]

This is the \emph{forward Dupire PDE} satisfied by call prices as a
function of maturity \(T\) and strike \(K\).

\paragraph{Remark.}
If interest rates or dividends are time-dependent, the formula must be
modified accordingly, but the essential idea remains the same:
local volatility is recovered from derivatives of call prices with
respect to maturity and strike.

\subsubsection{Interpretation}

Dupire's formula is remarkable because it expresses the local volatility
surface directly in terms of observable market quantities, namely the
prices of European call options across strikes and maturities.

\paragraph{What the formula says.}
The quantity
\[
\partial_T C(T,K)
\]
measures how the call price changes with maturity, while
\[
\partial_{KK}C(T,K)
\]
measures the curvature of the call price with respect to the strike.

Thus the local variance
\[
\sigma_{\mathrm{loc}}(T,K)^2
\]
is obtained by comparing:
\begin{itemize}[leftmargin=1.4em]
    \item the sensitivity of the option price to maturity,
    \item and the curvature of the option price with respect to strike.
\end{itemize}

\paragraph{Why this is important.}
The Black--Scholes model uses a single constant volatility parameter, so
it cannot reproduce a non-flat implied-volatility surface.
By contrast, Dupire's formula shows that, at least formally, one can
construct a deterministic function
\[
(T,K)\longmapsto \sigma_{\mathrm{loc}}(T,K)
\]
that reproduces exactly the observed market prices of European vanilla
options.

In this sense, the local volatility model is a model calibrated
\emph{directly} to the entire market surface of call prices.

\paragraph{What is local about it?}
The quantity \(\sigma_{\mathrm{loc}}(T,K)\) should be interpreted as the
instantaneous volatility assigned by the model when time is \(T\) and
the spot level is \(K\).
Thus the local volatility surface is read off from option prices, but it
is then reinserted into the diffusion dynamics of the underlying asset.

\paragraph{Key message.}
Dupire's formula provides a bridge between:
\begin{itemize}[leftmargin=1.4em]
    \item a \emph{cross-section of market option prices},
    \item and a \emph{state-dependent diffusion coefficient} for the
    underlying asset.
\end{itemize}

This is one of the main reasons why local volatility models are so
important in practice.

\subsubsection{Heuristic derivation from call prices}

We now explain heuristically where Dupire's formula comes from.

\paragraph{Step 1: write the call price in terms of the density.}
For simplicity, fix time \(0\), and denote by \(p(T,y)\) the density of
\(S_T\) under the risk-neutral measure.
Then the price of a European call with strike \(K\) and maturity \(T\)
is
\[
C(T,K)
=
e^{-rT}\int_0^\infty (y-K)^+\,p(T,y)\,\mathrm{d}y.
\]

\paragraph{Step 2: differentiate with respect to the strike.}
Differentiating once with respect to \(K\), we obtain
\[
\partial_K C(T,K)
=
-e^{-rT}\int_K^\infty p(T,y)\,\mathrm{d}y.
\]
Differentiating a second time gives
\[
\partial_{KK} C(T,K)
=
e^{-rT}p(T,K).
\]

Thus the second derivative of the call price with respect to the strike
recovers the risk-neutral density of the stock price.

\paragraph{Step 3: differentiate with respect to maturity.}
On the other hand,
\[
\partial_T C(T,K)
=
-rC(T,K)
+
e^{-rT}\int_0^\infty (y-K)^+\,\partial_T p(T,y)\,\mathrm{d}y.
\]

The density \(p(T,y)\) satisfies the forward Kolmogorov equation
\[
\partial_T p(T,y)
=
-\partial_y\big(ry\,p(T,y)\big)
+
\frac12\partial_{yy}\Big(\sigma_{\mathrm{loc}}(T,y)^2y^2p(T,y)\Big).
\]

Substituting this into the formula for \(\partial_T C(T,K)\), and then
integrating by parts, one obtains after simplification
\[
\partial_T C(T,K)
=
\frac12 \sigma_{\mathrm{loc}}(T,K)^2K^2\,\partial_{KK}C(T,K)
-rK\,\partial_K C(T,K).
\]

Rearranging yields Dupire's formula
\[
\sigma_{\mathrm{loc}}(T,K)^2
=
\frac{
\partial_T C(T,K)+rK\,\partial_K C(T,K)
}{
\frac12 K^2\,\partial_{KK}C(T,K)
}.
\]

\paragraph{Why this derivation is only heuristic.}
To make the argument fully rigorous, one must justify:
\begin{itemize}[leftmargin=1.4em]
    \item differentiation under the integral sign,
    \item the existence and smoothness of the density,
    \item and the integrations by parts.
\end{itemize}

For the purposes of this course, the key point is the structure of the
argument:
\begin{quote}
call prices determine the density through strike derivatives, and the
forward Kolmogorov equation then links this density to the local
volatility function.
\end{quote}

\subsection{Calibration to market option prices}

The main practical motivation for local volatility models is calibration
to the observed market surface of European option prices.

\paragraph{Calibration problem.}
Suppose that, for each maturity \(T\) and strike \(K\), the market
provides the call price
\[
C^{\mathrm{mkt}}(T,K),
\]
or equivalently the implied volatility surface
\[
\sigma_{\mathrm{imp}}(T,K).
\]
Since the Black--Scholes formula establishes a one-to-one relation
between option prices and implied volatilities, calibrating to market
option prices is equivalent to calibrating to the implied volatility
surface.  

The objective is therefore to choose a local volatility function
\[
\sigma_{\mathrm{loc}}(t,s)
\]
such that the model reproduces the observed market prices:
\[
C(0,S_0;T,K)=C^{\mathrm{mkt}}(T,K)
\qquad\text{for all quoted }(T,K).
\]

\paragraph{Why Black--Scholes cannot do this.}
The Black--Scholes model has only one volatility parameter \(\sigma\).
Hence it cannot fit an entire surface of option prices indexed by both
maturity and strike.
This mismatch is exactly what motivates the introduction of local
volatility: one replaces a single parameter by a full deterministic
function of time and spot.  

\paragraph{First naive idea.}
Starting from the backward pricing PDE,
\[
\partial_t C
+
rs\,\partial_s C
+
\frac12 \sigma_{\mathrm{loc}}(t,s)^2 s^2\,\partial_{ss}C
-rC=0,
\]
one might be tempted to solve for the local volatility as
\[
\sigma_{\mathrm{loc}}(t,s)^2
=
\frac{
rC-\partial_t C-rs\,\partial_s C
}{
\frac12 s^2\,\partial_{ss}C
}.
\]
However, this is not directly useful in practice, because these
derivatives are with respect to the current spot variable \(s\), while
market quotes are observed as functions of strike and maturity. This is
precisely why Dupire's forward formulation is needed.  

\paragraph{Calibration through Dupire's formula.}
Dupire's formula solves this problem by expressing the local volatility
directly in terms of derivatives of market call prices with respect to
maturity and strike:
\[
\sigma_{\mathrm{loc}}(T,K)^2
=
\frac{
\partial_T C(T,K)+rK\,\partial_K C(T,K)
}{
\frac12 K^2\,\partial_{KK}C(T,K)
}.
\]
Thus, at least formally, once the full call-price surface is known, the
local volatility surface can be recovered pointwise.

\paragraph{Interpolation issue.}
In practice, market quotes are only available on a finite and often
rather sparse grid of maturities and strikes.
Hence one cannot differentiate the market surface directly.
A first step is therefore to interpolate the observed data.

A common approach is the following:
\begin{itemize}[leftmargin=1.4em]
    \item for a fixed maturity, interpolate option prices across strikes,
    for instance with cubic splines;
    \item for a fixed strike, interpolate along the maturity direction,
    often in implied-volatility or total-variance space.  
\end{itemize}

In particular, if
\[
w(T,K):=\sigma_{\mathrm{BS}}(T,K)^2\,T
\]
denotes the total implied variance, then interpolating linearly in
\(w\) is a common and practical choice. 

\paragraph{Numerical differentiation and stability.}
Once an interpolated surface is available, one can compute numerical
derivatives in \(T\) and \(K\), and then apply Dupire's formula.

However, this step is numerically delicate:
\begin{itemize}[leftmargin=1.4em]
    \item the market surface may be noisy,
    \item numerical differentiation amplifies interpolation errors,
    \item for deep ITM or OTM options, the second derivative
    \(\partial_{KK}C(T,K)\) may become very small, which can cause
    instability in the Dupire formula. 
\end{itemize}

For this reason, practitioners often work in log-moneyness coordinates
and rewrite Dupire's formula in terms of total implied variance
\(w(T,K)\), which tends to be numerically more stable. 

\paragraph{Alternative calibration approach.}
Another possibility is to calibrate the local volatility surface by
optimization rather than pointwise inversion.
For a fixed maturity \(T\), one may search for a function
\[
\sigma_T(\cdot)=\sigma_{\mathrm{loc}}(T,\cdot)
\]
that minimizes the discrepancy between model prices and observed market
prices, for example in a least-squares sense:
\[
\min_{\sigma_T}
\sum_{j=1}^N
\big(
C_j^{\mathrm{obs}}-C(T,K_j;\sigma_T)
\big)^2.
\]
To avoid unstable or irregular surfaces, one often adds a regularization
term penalizing excessive curvature:
\[
\min_{\sigma_T}
\sum_{j=1}^N
\big(
C_j^{\mathrm{obs}}-C(T,K_j;\sigma_T)
\big)^2
+
\lambda\int_0^\infty
\big(\partial_{KK}\sigma_T(K)\big)^2\,\mathrm{d}K.
\]
This leads to a trade-off between exact fit and smoothness of the local
volatility surface.

\paragraph{Main lesson.}
Calibration is the central practical use of the local volatility model:
the whole purpose of introducing \(\sigma_{\mathrm{loc}}(t,s)\) is to
replace the single Black--Scholes volatility parameter by a function
flexible enough to fit the market cross-section of European option
prices.

\subsection{Strengths and limitations of local volatility}

The local volatility model is one of the most important extensions of
Black--Scholes, because it preserves much of its tractability while
allowing a much richer fit to market option prices.

\paragraph{Strengths.}

The first major strength of the model is that it can reproduce, at least
formally, the full market surface of European vanilla option prices at a
given date.
This is precisely the content of Dupire's formula.

The second advantage is that the model remains one-factor and Markovian.
The asset dynamics still have the form
\[
\mathrm{d}S_t
=
rS_t\,\mathrm{d}t
+
\sigma_{\mathrm{loc}}(t,S_t)S_t\,\mathrm{d}W_t,
\]
so pricing continues to rely on a single diffusion and a PDE very
similar to the Black--Scholes equation.

The third advantage is conceptual simplicity.
The model keeps the idea that all randomness comes from one Brownian
motion, and only modifies the local amplitude of the noise.
This makes it relatively easy to implement numerically once the local
volatility surface has been calibrated.

\paragraph{Practical appeal.}
Because of these features, local volatility models are widely used in
practice as arbitrage-free interpolation devices for the implied
volatility surface.
They provide a consistent way to price vanilla claims across a large set
of strikes and maturities.

\paragraph{Limitations.}

Despite these strengths, the local volatility model has important
limitations.

First, the model is entirely driven by a deterministic volatility
surface.
Once the current time and current spot are known, the instantaneous
volatility is fixed.
Hence the model does not introduce an independent volatility factor.
It therefore cannot capture the genuine randomness of volatility itself,
which is one of the main stylized facts of financial markets.  
Second, although the model fits the cross-section of vanilla option
prices at one date, it does not necessarily reproduce realistic
\emph{dynamics} of the implied volatility surface over time.
In particular, smile dynamics generated by local volatility models are
often considered less realistic than those produced by stochastic
volatility models.

Third, calibration is numerically fragile.
As already discussed, Dupire's formula requires differentiating market
option prices with respect to strike and maturity, and these numerical
derivatives can be unstable, especially in sparse or noisy markets.  

Fourth, the model may generate local volatility values that are
sensitive to interpolation choices and to small inconsistencies in the
market data.
In practice, extra smoothing and regularization are often needed.

\paragraph{Conceptual limitation.}
A useful way to summarize the situation is the following:
\begin{quote}
local volatility models are designed to fit today's smile, but not
necessarily tomorrow's smile dynamics.
\end{quote}

In other words, they are excellent \emph{static} calibration tools for
vanilla options, but they are less satisfactory when one wants a model
with realistic time-series behavior of volatility.

\paragraph{Comparison with stochastic volatility.}
This is precisely why stochastic volatility models are introduced.
Whereas local volatility explains the smile by making volatility a
deterministic function of \((t,S_t)\), stochastic volatility models
introduce a separate random factor for volatility itself.
They typically provide a better dynamic description of market smiles and
skews, at the cost of a more complex model.

\paragraph{Main lesson.}
The local volatility model should therefore be viewed as a very useful
intermediate model:
\begin{itemize}[leftmargin=1.4em]
    \item richer and more realistic than Black--Scholes,
    \item still tractable and calibratable to vanilla options,
    \item but less realistic than stochastic volatility models when it
    comes to volatility dynamics.
\end{itemize}

% ============================================================
\section{Stochastic volatility models}
% ============================================================

\subsection{Why local volatility is not enough}

The local volatility model is a major improvement over Black--Scholes,
because it allows one to fit the whole cross-section of European vanilla
option prices observed at a given date.
However, despite this important success, it is still not a fully
satisfactory model of market dynamics.

\paragraph{A first limitation: volatility is not genuinely random.}
In a local volatility model, the asset price satisfies
\[
\mathrm{d}S_t
=
rS_t\,\mathrm{d}t
+
\sigma_{\mathrm{loc}}(t,S_t)S_t\,\mathrm{d}W_t.
\]
Thus, once the current time \(t\) and the current spot level \(S_t\) are
known, the instantaneous volatility is completely determined.
The model remains one-factor and Markovian, and there is no additional
source of randomness driving volatility itself.

This is at odds with market evidence.
In practice, volatility behaves like a random state variable:
it clusters through time, varies in persistent ways, and reacts
strongly to market movements.
A deterministic function \(\sigma_{\mathrm{loc}}(t,s)\) cannot reproduce
this feature.

\paragraph{A second limitation: unrealistic smile dynamics.}
The local volatility model is calibrated so as to match the implied
volatility surface at the current date.
In this sense, it is a very good \emph{static} model.

However, once the model is evolved forward in time, the future
dynamics of the smile are entirely determined by the deterministic
function \(\sigma_{\mathrm{loc}}(t,s)\).
As a consequence, the way the implied volatility surface moves when the
spot moves is often unrealistic.

In particular, local volatility models typically produce smile dynamics
that are too strongly tied to the spot level.
Practitioners often summarize this by saying that local volatility
models fit today's smile, but do not generate realistic future smile
movements.

\paragraph{A third limitation: poor description of volatility time series.}
Because the model has only one Brownian driver, all randomness is carried
by the asset price itself.
There is no separate process representing the volatility level.

As a consequence, the model does not capture important empirical
features such as:
\begin{itemize}[leftmargin=1.4em]
    \item volatility clustering,
    \item persistence of high- and low-volatility regimes,
    \item a separate source of uncertainty attached to future volatility.
\end{itemize}

This is a serious limitation for products whose value depends not only
on the terminal distribution of the asset, but also on the future
evolution of volatility.

\paragraph{A fourth limitation: sensitivity to calibration.}
The local volatility surface is recovered from market option prices
through Dupire-type formulas or related calibration procedures.
In practice, this is numerically delicate:
\begin{itemize}[leftmargin=1.4em]
    \item market quotes are observed only on a finite grid,
    \item interpolation is required,
    \item numerical differentiation may amplify noise,
    \item the resulting local volatility surface may be unstable.
\end{itemize}

Thus, although the model is elegant in theory, its practical
implementation can be fragile.

\paragraph{A useful summary.}
The local volatility model should be viewed as a model that reproduces
the \emph{marginal law} of the asset well enough to fit vanilla option
prices at one date, but not as a model that reproduces correctly the
joint dynamics of spot and volatility through time.

\paragraph{Why this motivates stochastic volatility.}
These limitations naturally lead to the introduction of stochastic
volatility models.

Instead of imposing that volatility is a deterministic function of
\((t,S_t)\), one introduces a new random factor \(V_t\), and writes a
system of the form
\[
\mathrm{d}S_t
=
rS_t\,\mathrm{d}t
+
\sqrt{V_t}\,S_t\,\mathrm{d}W_t^{(1)},
\]
together with a second equation describing the dynamics of \(V_t\).

This allows volatility itself to fluctuate randomly, and leads to richer
and more realistic smile dynamics.

\paragraph{Main lesson.}
Local volatility is an excellent model for static calibration to vanilla
options, but it is not sufficient when one wants:
\begin{itemize}[leftmargin=1.4em]
    \item a realistic dynamic model of the implied volatility surface,
    \item a separate source of volatility risk,
    \item or a better description of the time-series behavior of
    financial markets.
\end{itemize}

This is the main reason why stochastic volatility models are introduced.

\subsection{General structure of a stochastic volatility model}

We now move from local volatility models to a broader class of models in
which volatility itself is random.

\paragraph{Basic idea.}
In the Black--Scholes model, volatility is constant.
In a local volatility model, volatility is a deterministic function of
time and spot,
\[
\sigma_{\mathrm{loc}}(t,S_t).
\]
In both cases, once the current state of the stock is known, the
instantaneous volatility is fully determined.

A stochastic volatility model goes one step further:
it introduces an additional random factor, usually denoted by \(V_t\) or
\(\nu_t\), which drives the volatility of the asset.

\paragraph{Generic form of the model.}
Under the risk-neutral measure, a stochastic volatility model typically
takes the form
\[
\mathrm{d}S_t
=
rS_t\,\mathrm{d}t
+
\sigma(t,S_t,V_t)\,S_t\,\mathrm{d}W_t^{(1)},
\]
together with a second equation describing the volatility factor:
\[
\mathrm{d}V_t
=
b(t,V_t)\,\mathrm{d}t
+
a(t,V_t)\,\mathrm{d}W_t^{(2)}.
\]

The two Brownian motions are generally allowed to be correlated:
\[
\mathrm{d}\langle W^{(1)},W^{(2)}\rangle_t=\rho\,\mathrm{d}t,
\qquad -1\le \rho \le 1.
\]

Thus the model is usually a two-factor diffusion:
\begin{itemize}[leftmargin=1.4em]
    \item one factor for the asset price,
    \item one factor for the volatility level.
\end{itemize}

\paragraph{Interpretation of the volatility factor.}
The process \(V_t\) may represent:
\begin{itemize}[leftmargin=1.4em]
    \item the instantaneous variance,
    \item the instantaneous volatility,
    \item or a latent factor driving volatility.
\end{itemize}

A very common convention is to interpret \(V_t\) as the instantaneous
variance and to write
\[
\mathrm{d}S_t
=
rS_t\,\mathrm{d}t
+
\sqrt{V_t}\,S_t\,\mathrm{d}W_t^{(1)}.
\]

In that case, \(V_t\) must remain non-negative.

\paragraph{Why correlation matters.}
The correlation parameter \(\rho\) plays an important role in practice.
It couples the shocks on the asset price and on volatility.
In particular, negative values of \(\rho\) are often used to reproduce
the equity volatility skew:
when the asset price drops, volatility tends to increase.

Thus stochastic volatility models can naturally generate asymmetric
smiles and skews.

\paragraph{Comparison with local volatility.}
The key difference with local volatility is that the instantaneous
volatility is no longer determined solely by \((t,S_t)\).
Even if the current stock price \(S_t\) is known, the future evolution
of volatility still depends on the additional random factor \(V_t\).

Therefore, stochastic volatility models are no longer one-factor models:
they introduce an additional source of risk.

\paragraph{Markovian structure.}
Although the stock price process alone is no longer Markovian in
general, the pair
\[
(S_t,V_t)
\]
is typically Markovian.
As a consequence, derivative prices are naturally represented as
functions of two state variables:
\[
u(t,s,v).
\]

This leads to pricing PDEs in higher dimension than in Black--Scholes or
local volatility models.

\paragraph{Pricing point of view.}
For a European payoff \(g(S_T)\), the risk-neutral pricing formula
becomes
\[
u(t,s,v)
=
\E^\Qbb\!\left[
e^{-r(T-t)}g(S_T)
\;\middle|\;
S_t=s,\ V_t=v
\right].
\]

Thus the price depends not only on the current level of the asset, but
also on the current level of volatility.

\paragraph{Main advantage.}
The main benefit of stochastic volatility models is that they provide a
more realistic description of market dynamics:
\begin{itemize}[leftmargin=1.4em]
    \item volatility is genuinely random,
    \item the model can capture volatility clustering,
    \item smile and skew dynamics are richer,
    \item the model is better suited to products sensitive to future
    volatility.
\end{itemize}

\paragraph{Main drawback.}
The price to pay is increased complexity:
\begin{itemize}[leftmargin=1.4em]
    \item the model is higher-dimensional,
    \item calibration is more delicate,
    \item pricing and hedging are more involved,
    \item closed-form formulas are rare, except in a few special models.
\end{itemize}

\paragraph{Examples.}
Several well-known models belong to this class:
\begin{itemize}[leftmargin=1.4em]
    \item the Heston model,
    \item the Hull--White stochastic volatility model,
    \item SABR-type models in interest-rate and FX markets.
\end{itemize}

In the next subsection, we focus on the Heston model, which is one of
the most important and widely used stochastic volatility models.


\subsection{Example: a two-factor stochastic volatility model}

We now present a simple example of a stochastic volatility model with
two state variables: the asset price and the volatility factor.

\paragraph{Model specification.}
Under the risk-neutral measure, consider the system
\[
\mathrm{d}S_t
=
rS_t\,\mathrm{d}t
+
\sqrt{V_t}\,S_t\,\mathrm{d}W_t^{(1)},
\]
\[
\mathrm{d}V_t
=
\kappa(\theta-V_t)\,\mathrm{d}t
+
\eta\sqrt{V_t}\,\mathrm{d}W_t^{(2)},
\]
where
\[
\mathrm{d}\langle W^{(1)},W^{(2)}\rangle_t=\rho\,\mathrm{d}t,
\qquad -1\le \rho\le 1.
\]

Here:
\begin{itemize}[leftmargin=1.4em]
    \item \(S_t\) is the asset price,
    \item \(V_t\) is the instantaneous variance,
    \item \(r\) is the risk-free rate,
    \item \(\kappa>0\) is the mean-reversion speed of the variance,
    \item \(\theta>0\) is the long-run variance level,
    \item \(\eta>0\) is the volatility of volatility,
    \item \(\rho\) is the correlation between the shocks on the asset and on the variance.
\end{itemize}

\paragraph{Why this is a two-factor model.}
The model is driven by two sources of randomness:
\begin{itemize}[leftmargin=1.4em]
    \item one Brownian motion \(W^{(1)}\) acting directly on the stock price,
    \item one Brownian motion \(W^{(2)}\) acting on the variance process.
\end{itemize}
Thus the state variable is the pair
\[
(S_t,V_t),
\]
and the model is genuinely two-dimensional.

\paragraph{Interpretation of the variance dynamics.}
The variance process \(V_t\) is mean reverting.
Indeed, its drift is
\[
\kappa(\theta-V_t).
\]
Hence:
\begin{itemize}[leftmargin=1.4em]
    \item if \(V_t>\theta\), the drift is negative and tends to bring the variance down,
    \item if \(V_t<\theta\), the drift is positive and tends to push the variance up.
\end{itemize}

This is consistent with the empirical observation that volatility tends
to fluctuate around a long-run level rather than drift away indefinitely.

\paragraph{Role of the volatility of volatility.}
The coefficient \(\eta\) measures the amplitude of the random
fluctuations of the variance process.
A larger value of \(\eta\) means that volatility itself is more random.

This parameter has a strong impact on option prices, especially on the
shape of the implied volatility smile.

\paragraph{Role of the correlation.}
The correlation \(\rho\) is one of the most important parameters of the
model.
It controls the interaction between the stock and the variance.

In equity markets, one often uses \(\rho<0\).
This reflects the fact that downward moves in the stock are frequently
associated with increases in volatility.
Such a negative correlation naturally generates a downward volatility
skew.

\paragraph{Log-price dynamics.}
Applying It\^o's formula to \(\log S_t\), we obtain
\[
\mathrm{d}\log S_t
=
\Big(r-\frac12 V_t\Big)\mathrm{d}t
+
\sqrt{V_t}\,\mathrm{d}W_t^{(1)}.
\]

Compared with Black--Scholes, the key difference is that the
instantaneous variance is no longer constant, but equal to the random
process \(V_t\).

\paragraph{Pricing point of view.}
If \(g(S_T)\) is the payoff of a European derivative, then its price is
naturally written as
\[
u(t,s,v)
=
\E^\Qbb\!\left[
e^{-r(T-t)}g(S_T)
\;\middle|\;
S_t=s,\ V_t=v
\right].
\]

Thus the option price depends on two state variables:
\begin{itemize}[leftmargin=1.4em]
    \item the current stock price \(s\),
    \item the current variance level \(v\).
\end{itemize}

\paragraph{Associated PDE.}
The infinitesimal generator of the two-dimensional diffusion
\((S_t,V_t)\) is
\[
\mathcal L u
=
rs\,\partial_s u
+
\kappa(\theta-v)\,\partial_v u
+
\frac12 vs^2\,\partial_{ss}u
+
\rho\eta vs\,\partial_{sv}u
+
\frac12\eta^2 v\,\partial_{vv}u.
\]

Therefore the option price satisfies the PDE
\[
\partial_t u+\mathcal L u-r u=0,
\qquad
u(T,s,v)=g(s).
\]

This illustrates one of the main differences with local volatility
models: the pricing equation is now two-dimensional.

\paragraph{What this model captures.}
This type of model is able to reproduce several important features of
financial markets:
\begin{itemize}[leftmargin=1.4em]
    \item stochastic fluctuations of volatility,
    \item volatility clustering,
    \item skewed implied volatility surfaces,
    \item richer smile dynamics than local volatility models.
\end{itemize}

\paragraph{Connection with the Heston model.}
The model above is in fact the classical Heston-type specification.
It is one of the most important stochastic volatility models in
practice, because it is rich enough to produce realistic smile effects
while remaining analytically tractable in some cases.

In the next subsection, we discuss this model in more detail and explain
why it plays such a central role in option pricing.

\subsection{The Heston model}

Among stochastic volatility models, the Heston model is one of the most
important and widely used in practice.
It combines a relatively rich volatility dynamics with a level of
tractability that makes it especially attractive for pricing and
calibration.

\subsubsection{Dynamics}

Under the risk-neutral measure, the Heston model is given by
\[
\mathrm{d}S_t
=
rS_t\,\mathrm{d}t
+
\sqrt{V_t}\,S_t\,\mathrm{d}W_t^{(1)},
\]
\[
\mathrm{d}V_t
=
\kappa(\theta-V_t)\,\mathrm{d}t
+
\eta\sqrt{V_t}\,\mathrm{d}W_t^{(2)},
\]
with
\[
\mathrm{d}\langle W^{(1)},W^{(2)}\rangle_t=\rho\,\mathrm{d}t,
\qquad -1\le \rho\le 1.
\]

Here:
\begin{itemize}[leftmargin=1.4em]
    \item \(S_t\) is the stock price,
    \item \(V_t\) is the instantaneous variance,
    \item \(r\) is the risk-free rate,
    \item \(\kappa>0\) is the mean-reversion speed,
    \item \(\theta>0\) is the long-run variance level,
    \item \(\eta>0\) is the volatility of volatility,
    \item \(\rho\) is the correlation between the stock and variance shocks.
\end{itemize}

\paragraph{Interpretation of the parameters.}

The parameter \(\kappa\) controls how quickly the variance tends to
return to its long-run level \(\theta\).
If \(\kappa\) is large, the variance reverts rapidly; if \(\kappa\) is
small, the variance remains persistent for longer periods.

The parameter \(\eta\) controls the randomness of volatility itself.
Larger values of \(\eta\) produce stronger fluctuations in the variance
process and therefore more pronounced smile effects.

The parameter \(\theta\) determines the average variance level around
which the process fluctuates.

\paragraph{Positivity of the variance.}

Since \(V_t\) represents an instantaneous variance, it should remain
non-negative.
The square-root diffusion structure is designed precisely for this
purpose.
A sufficient condition ensuring strict positivity is the Feller
condition
\[
2\kappa\theta \ge \eta^2.
\]
Even when this condition is not imposed, the model is still widely used
in practice, although the behavior of the variance process near zero
must then be treated with care.

\paragraph{Log-price dynamics.}

Applying It\^o's formula to \(\log S_t\), we obtain
\[
\mathrm{d}\log S_t
=
\Big(r-\frac12V_t\Big)\mathrm{d}t
+
\sqrt{V_t}\,\mathrm{d}W_t^{(1)}.
\]

Thus the instantaneous variance of log-returns is exactly \(V_t\), which
is now random and time-varying.

\paragraph{Pricing point of view.}

If \(g(S_T)\) is the payoff of a European derivative, its price is
written as
\[
u(t,s,v)
=
\E^\Qbb\!\left[
e^{-r(T-t)}g(S_T)
\;\middle|\;
S_t=s,\ V_t=v
\right].
\]

As in the previous subsection, the price depends on two state variables,
and the associated pricing PDE is two-dimensional.

\subsubsection{Role of correlation}

The correlation parameter \(\rho\) plays a crucial role in the Heston
model.

\paragraph{Why correlation matters.}

The Brownian motion \(W^{(1)}\) drives the stock price, while
\(W^{(2)}\) drives the variance process.
If these two Brownian motions were independent, then shocks on the stock
and shocks on volatility would be unrelated.
In practice, this is not realistic, especially in equity markets.

Indeed, empirical evidence suggests that large downward moves in stock
prices are often associated with increases in volatility.
This phenomenon is sometimes referred to as the \emph{leverage effect}.

\paragraph{Negative correlation.}

When \(\rho<0\), a negative shock on the stock tends to be accompanied by
a positive shock on variance.
Thus falling prices are associated with rising volatility.

This is exactly the mechanism needed to generate the downward skew
observed in equity implied volatility surfaces.

\paragraph{Positive correlation.}

If \(\rho>0\), upward moves in the stock tend to be accompanied by
increases in volatility.
This is less typical in equity markets, but may be relevant in other
asset classes.

\paragraph{Special case \(\rho=0\).}

If \(\rho=0\), stock and variance shocks are uncorrelated.
The model can still generate a volatility smile because volatility is
random, but it loses the strong asymmetry needed to reproduce a
pronounced skew.

\paragraph{Economic interpretation.}

The correlation parameter is therefore not just a technical feature of
the model.
It is one of the main mechanisms through which the Heston model produces
asymmetric option prices and realistic implied-volatility patterns.

\subsubsection{Qualitative effect on the smile/skew}

The Heston model is particularly valued because it can generate implied
volatility smiles and skews in a natural way.

\paragraph{Effect of stochastic volatility itself.}

Even if \(\rho=0\), the randomness of the variance process already makes
the distribution of returns richer than in Black--Scholes.
In particular, the model can generate heavier tails and more curvature
in implied volatility across strikes.

Thus stochastic volatility alone already tends to produce a smile.

\paragraph{Effect of negative correlation.}

When \(\rho<0\), the model produces a downward skew in implied
volatility:
low strikes tend to have higher implied volatilities than high strikes.

This is the pattern typically observed in equity markets.
The reason is that large downward stock moves are associated with higher
future volatility, which increases the value of out-of-the-money puts.

\paragraph{Effect of \(\eta\).}

The parameter \(\eta\), often called the volatility of volatility, has a
strong impact on the curvature of the smile.
A larger \(\eta\) produces more randomness in volatility and therefore a
more pronounced smile or skew.

\paragraph{Effect of \(\kappa\) and \(\theta\).}

The long-run variance level \(\theta\) influences the overall level of
implied volatility, while the mean-reversion speed \(\kappa\) affects
how persistent volatility shocks are.
Together, these parameters influence both the level and the term
structure of implied volatility.

\paragraph{Comparison with local volatility.}

Unlike local volatility models, the Heston model does not merely fit a
static implied volatility surface through a deterministic function of
time and spot.
Instead, it introduces a genuine random volatility factor, which leads
to more realistic smile dynamics over time.

This is one of the main reasons why stochastic volatility models are
preferred when one wants not only a fit to today's market prices, but
also a plausible evolution of future implied volatility surfaces.

\paragraph{Main lesson.}

The Heston model combines three important features:
\begin{itemize}[leftmargin=1.4em]
    \item random volatility,
    \item mean reversion of variance,
    \item and correlation between spot and volatility shocks.
\end{itemize}

Together, these ingredients allow the model to generate realistic
smiles and skews while remaining sufficiently tractable for pricing and
calibration.

\subsection{Pricing and hedging issues under stochastic volatility}

Stochastic volatility models are appealing because they provide a much
richer description of market dynamics than Black--Scholes or local
volatility models.
However, this additional realism comes at a price: both pricing and
hedging become more difficult.

\paragraph{Pricing is no longer one-dimensional.}

In the Black--Scholes model, the option price depends only on time and
the current stock price:
\[
u=u(t,s).
\]
In a stochastic volatility model, the relevant state variable is
typically the pair
\[
(S_t,V_t),
\]
so the option price becomes a function
\[
u=u(t,s,v).
\]

As a result, the pricing PDE is no longer one-dimensional, but
two-dimensional.
For example, in a Heston-type model, the pricing equation takes the form
\[
\partial_t u+\mathcal L u-r u=0,
\qquad
u(T,s,v)=g(s),
\]
where \(\mathcal L\) is the generator of the two-factor diffusion.
This already makes numerical pricing more expensive than in the
Black--Scholes or local volatility settings.

\paragraph{Closed-form formulas are rare.}

In Black--Scholes, European call and put prices are given by explicit
closed-form formulas.
In stochastic volatility models, such explicit expressions are usually
not available.

Some models, such as the Heston model, remain partially tractable and
allow semi-explicit pricing formulas based on Fourier methods or
characteristic functions.
But in general, one often has to rely on:
\begin{itemize}[leftmargin=1.4em]
    \item finite-difference methods,
    \item Monte Carlo simulation,
    \item Fourier techniques,
    \item or approximation formulas.
\end{itemize}

Thus stochastic volatility models are significantly more demanding from
a numerical point of view.

\paragraph{A new source of risk appears.}

A major conceptual difference with Black--Scholes is that volatility is
now random.
Therefore, in addition to the uncertainty coming from the stock price,
there is an extra source of randomness coming from the volatility factor.

In a model of the form
\[
\mathrm{d}S_t
=
rS_t\,\mathrm{d}t
+
\sqrt{V_t}\,S_t\,\mathrm{d}W_t^{(1)},
\qquad
\mathrm{d}V_t
=
b(t,V_t)\,\mathrm{d}t
+
a(t,V_t)\,\mathrm{d}W_t^{(2)},
\]
there are typically two Brownian drivers.
Hence the market is no longer complete if one trades only in the stock
and the money market account.

This means that one cannot, in general, perfectly replicate every
derivative written on the stock.

\paragraph{Consequence for hedging.}

In the Black--Scholes model, delta hedging is enough to eliminate all
risk, because the model has only one source of randomness and one risky
asset.
Under stochastic volatility, delta hedging is no longer sufficient.

Indeed, a portfolio hedged only with respect to movements in the stock
still remains exposed to changes in the volatility factor.
This is often called \emph{volatility risk}.

As a consequence, hedging requires additional instruments, typically
other liquid options, in order to control sensitivity not only to the
spot, but also to the volatility factor.

\paragraph{Greeks and volatility exposure.}

In stochastic volatility models, the usual delta and gamma remain
important, but other sensitivities become central.
In particular:
\begin{itemize}[leftmargin=1.4em]
    \item \emph{vega} measures sensitivity to changes in volatility,
    \item the sensitivity to the variance factor \(V_t\) itself becomes
    a key object,
    \item correlation risk may also matter through the parameter \(\rho\).
\end{itemize}

Thus risk management becomes multidimensional:
one no longer hedges only the spot risk, but also part of the volatility
risk.

\paragraph{Model incompleteness.}

Because stochastic volatility introduces an additional non-traded source
of risk, the model is usually incomplete.
This has several consequences:
\begin{itemize}[leftmargin=1.4em]
    \item perfect replication is generally impossible,
    \item option prices depend on the choice of risk-neutral dynamics for
    volatility,
    \item hedging becomes approximate rather than exact.
\end{itemize}

In practice, one therefore works with calibrated models and hedging
strategies that are designed to perform well empirically, rather than
with exact replication arguments as in Black--Scholes.

\paragraph{Numerical challenges.}

Stochastic volatility models also raise several numerical issues.

First, the pricing PDE is higher-dimensional, which makes finite
difference methods more expensive.

Second, Monte Carlo methods remain available, but they require the
simulation of both the stock and the volatility process.

Third, particular care is needed when the variance process must remain
non-negative, as in the Heston model.

Fourth, calibration itself is more involved, since several parameters
must be estimated simultaneously:
\[
\kappa,\qquad \theta,\qquad \eta,\qquad \rho,\qquad V_0.
\]

\paragraph{Practical trade-off.}

The advantage of stochastic volatility models is that they reproduce
market smiles and skews much more realistically than Black--Scholes or
local volatility.
The drawback is that they are harder to analyze, harder to calibrate,
and harder to hedge.

Thus they represent a trade-off:
\begin{itemize}[leftmargin=1.4em]
    \item more realism,
    \item but less tractability.
\end{itemize}

\paragraph{Main lesson.}

Under stochastic volatility, pricing and hedging become intrinsically
more complex because volatility itself is random.
This enriches the model considerably, but it also destroys the exact
replication paradigm of Black--Scholes and replaces it with a more
realistic, but more difficult, framework.

\subsection{Calibration and practical interpretation}

One of the main reasons stochastic volatility models are so important in
practice is that they provide a flexible framework for fitting market
option prices while retaining a realistic interpretation of volatility
risk.

\paragraph{Why calibration is needed.}

A stochastic volatility model contains several parameters, for instance
in the Heston model
\[
\kappa,\qquad \theta,\qquad \eta,\qquad \rho,\qquad V_0.
\]
These parameters are not directly observable in the market.
They must therefore be chosen so that the model reproduces, as closely
as possible, the prices of liquidly traded options.

In practice, the calibration target is not usually the raw option-price
surface itself, but rather the implied volatility surface, since this is
how market option prices are commonly quoted.

\paragraph{Calibration objective.}

Suppose that the market provides option prices or implied volatilities
for a set of strikes and maturities
\[
(K_i,T_j).
\]
The goal is to choose the model parameters so that the model prices
\[
C^{\mathrm{model}}(K_i,T_j)
\]
match the observed market prices
\[
C^{\mathrm{mkt}}(K_i,T_j),
\]
or equivalently that the model implied volatilities match the market
implied volatilities.

A typical calibration problem therefore takes the form
\[
\min_{\Theta}
\sum_{i,j}
w_{ij}
\Big(
C^{\mathrm{model}}(K_i,T_j;\Theta)
-
C^{\mathrm{mkt}}(K_i,T_j)
\Big)^2,
\]
or, more commonly in practice,
\[
\min_{\Theta}
\sum_{i,j}
w_{ij}
\Big(
\sigma_{\mathrm{imp}}^{\mathrm{model}}(K_i,T_j;\Theta)
-
\sigma_{\mathrm{imp}}^{\mathrm{mkt}}(K_i,T_j)
\Big)^2,
\]
where \(\Theta\) denotes the vector of model parameters and
\(w_{ij}\) are chosen weights.

\paragraph{Why implied volatilities are often used.}

There are several practical reasons for calibrating to implied
volatility rather than directly to prices:
\begin{itemize}[leftmargin=1.4em]
    \item implied volatility normalizes prices across strikes and maturities,
    \item market practitioners naturally think in terms of volatility
    surfaces,
    \item errors measured in implied volatility are often easier to
    interpret than raw price errors.
\end{itemize}

\paragraph{Typical role of the parameters.}

A useful feature of stochastic volatility models is that the parameters
have an economic and geometric interpretation.

In a Heston-type model:
\begin{itemize}[leftmargin=1.4em]
    \item \(V_0\) controls the short-maturity level of volatility,
    \item \(\theta\) controls the long-run level of variance,
    \item \(\kappa\) controls how quickly the variance reverts toward
    \(\theta\),
    \item \(\eta\) controls the amount of randomness in volatility and
    therefore the curvature of the smile,
    \item \(\rho\) controls the asymmetry of the smile, in particular the
    slope of the skew.
\end{itemize}

Thus calibration is not only a numerical procedure: it also gives an
economic interpretation of the observed implied volatility surface.

\paragraph{Practical difficulties.}

Calibration under stochastic volatility is significantly more delicate
than under local volatility.

First, model prices are usually more expensive to compute, since one
must solve a two-dimensional PDE, use Fourier methods, or rely on Monte
Carlo simulation.

Second, the objective function may have several local minima, so the
quality of calibration may depend on the initial guess and on the chosen
optimization routine.

Third, different parameter combinations may lead to similar implied
volatility surfaces, so calibration can be unstable or poorly
identified.

Fourth, one often wants not only a good fit at one date, but also
parameters that remain economically reasonable and stable through time.

\paragraph{Static fit versus dynamic interpretation.}

A key point is that calibration should not be judged only by how well
the model fits today's smile.

A model is also intended to describe how the volatility surface will
move in the future.
This is one of the main differences with local volatility models:
stochastic volatility models are typically less exact in their static
fit, but more realistic in their dynamics.

Hence, in practice, one often accepts a slightly less perfect fit at one
given date in exchange for a more plausible evolution of the smile over
time.

\paragraph{Practical interpretation.}

From a market perspective, stochastic volatility models provide a bridge
between two viewpoints:
\begin{itemize}[leftmargin=1.4em]
    \item a \emph{cross-sectional} viewpoint, through calibration to
    today's option surface,
    \item a \emph{dynamic} viewpoint, through the random evolution of the
    volatility factor.
\end{itemize}

This makes them particularly attractive for:
\begin{itemize}[leftmargin=1.4em]
    \item pricing exotic derivatives,
    \item managing vega and smile risk,
    \item stress testing volatility scenarios,
    \item and designing hedging strategies that remain meaningful over
    time.
\end{itemize}

\paragraph{Comparison with local volatility.}

It is useful to summarize the contrast as follows:
\begin{itemize}[leftmargin=1.4em]
    \item local volatility is designed to fit exactly the current
    vanilla surface,
    \item stochastic volatility is designed to provide a more realistic
    dynamic description of volatility risk.
\end{itemize}

Thus stochastic volatility models are often preferred when the
derivative being priced is sensitive not only to today's smile, but also
to its future evolution.

\paragraph{Main lesson.}

Calibration gives numerical values to the parameters of a stochastic
volatility model, but these parameters should not be viewed as purely
technical constants.
They describe:
\begin{itemize}[leftmargin=1.4em]
    \item the level of volatility,
    \item its persistence,
    \item its randomness,
    \item and its interaction with the underlying asset.
\end{itemize}

In this sense, calibration is both a pricing tool and a way of
translating market option prices into a model-based description of
volatility risk.

\subsection{Comparison with local volatility}

We now compare local volatility and stochastic volatility models.

Both approaches were introduced to go beyond the Black--Scholes model
and to account for the fact that market option prices exhibit smiles,
skews, and maturity effects.
However, they do so in very different ways.

\paragraph{Local volatility: deterministic volatility surface.}

In a local volatility model, the stock price satisfies
\[
\mathrm{d}S_t
=
rS_t\,\mathrm{d}t
+
\sigma_{\mathrm{loc}}(t,S_t)S_t\,\mathrm{d}W_t,
\]
where the volatility coefficient is a deterministic function of time and
spot.

Thus, once \(t\) and \(S_t\) are known, the instantaneous volatility is
completely determined.
The model remains one-factor and Markovian.

Its main strength is that it can be calibrated, at least formally, so as
to reproduce exactly the observed surface of European vanilla option
prices at a given date.

\paragraph{Stochastic volatility: random volatility factor.}

In a stochastic volatility model, the stock price depends on an
additional random factor \(V_t\), for example through
\[
\mathrm{d}S_t
=
rS_t\,\mathrm{d}t
+
\sqrt{V_t}\,S_t\,\mathrm{d}W_t^{(1)}.
\]
The volatility is therefore no longer a deterministic function of
\((t,S_t)\), but a genuinely random process.

As a consequence, the model is typically two-factor, with state variable
\[
(S_t,V_t),
\]
and the future behavior of volatility is itself uncertain.

\paragraph{Static fit versus dynamic realism.}

This is the main contrast between the two approaches.

Local volatility models are designed primarily for \emph{static}
calibration.
They fit today's implied volatility surface very well, but their future
smile dynamics are often considered unrealistic.

By contrast, stochastic volatility models are typically less exact in
their static fit, but they provide a more realistic \emph{dynamic}
description of market volatility.
They allow the implied volatility surface itself to evolve in a random
and economically meaningful way.

\paragraph{Number of risk factors.}

Another major difference lies in the number of sources of randomness.

In local volatility, the only source of randomness is the Brownian
motion driving the stock price.
The model is therefore one-factor.

In stochastic volatility, there is usually an additional Brownian motion
driving the volatility factor.
The model becomes multi-factor, and this creates additional sources of
risk, especially volatility risk.

\paragraph{Consequences for hedging.}

Because local volatility models remain one-factor, they preserve many of
the structural features of Black--Scholes.
In particular, they are still complete in the idealized setting, and
delta hedging remains the central replication tool.

By contrast, stochastic volatility models are generally incomplete.
Delta hedging is no longer sufficient, because the derivative is exposed
not only to spot risk, but also to volatility risk.
In practice, hedging often requires additional liquid options.

\paragraph{Consequences for computation.}

Local volatility models remain relatively tractable.
The pricing PDE is still one-dimensional:
\[
u=u(t,s),
\]
and the numerical methods are close in spirit to the Black--Scholes
case.

Stochastic volatility models are more demanding computationally.
The price now depends on two state variables,
\[
u=u(t,s,v),
\]
so the associated PDE is higher-dimensional, and numerical pricing
becomes more expensive.

\paragraph{Calibration point of view.}

Local volatility calibration is usually based directly on the market
surface of vanilla option prices, for instance through Dupire's formula.

Stochastic volatility calibration is typically formulated as an inverse
problem in a finite-dimensional parameter space.
This is often more stable conceptually, but the fit is usually less
exact than in local volatility.

\paragraph{A useful summary.}

A simple way to summarize the difference is the following:
\begin{quote}
local volatility is excellent for fitting today's smile, while
stochastic volatility is better for modeling tomorrow's smile.
\end{quote}

\paragraph{Which model should one prefer?}

The answer depends on the objective.

If the main goal is to match the current cross-section of vanilla option
prices as accurately as possible, then a local volatility model is very
attractive.

If the goal is to model volatility risk dynamically, to price exotic
products, or to obtain more realistic smile dynamics over time, then
stochastic volatility models are usually preferred.

\paragraph{Main lesson.}

The two approaches should not be viewed as competitors in a purely
absolute sense.
Rather, they answer different modeling needs:
\begin{itemize}[leftmargin=1.4em]
    \item local volatility emphasizes exact static calibration,
    \item stochastic volatility emphasizes realistic dynamic behavior.
\end{itemize}

In practice, both families of models are fundamental, and the choice
between them depends on the balance one wants between tractability,
calibration quality, and realism.

% ============================================================
\section{Stochastic interest-rate modeling}
% ============================================================

\subsection{Why constant interest rates are unrealistic}

In the Black--Scholes model, the interest rate is usually assumed to be
constant.
This assumption is very convenient mathematically, since it leads to a
simple discount factor
\[
e^{-r(T-t)},
\]
and allows one to price derivatives by discounting expected future
payoffs at a fixed rate.

However, this assumption is often too restrictive in practice.

\paragraph{Interest rates vary across time.}

In real markets, interest rates are not constant.
They fluctuate over time in response to monetary policy, inflation,
macroeconomic conditions, and market expectations.
As a consequence, the value of one unit of cash invested in the money
market account is itself uncertain over long horizons.

Thus the short rate should more realistically be modeled as a process
\[
r_t,
\]
rather than as a constant number \(r\).

\paragraph{The term structure is not flat.}

If the interest rate were constant, then the prices of zero-coupon bonds
would satisfy
\[
P(t,T)=e^{-r(T-t)},
\]
so the entire yield curve would be flat.

In reality, this is not what one observes.
For different maturities \(T\), zero-coupon bonds have different yields,
and the resulting curve may be:
\begin{itemize}[leftmargin=1.4em]
    \item upward sloping,
    \item downward sloping,
    \item or hump-shaped.
\end{itemize}

Hence one cannot describe the whole term structure of interest rates
with a single constant parameter.

\paragraph{Bond prices carry their own source of risk.}

Under constant rates, bond prices are deterministic functions of time.
In particular, the discount factor is known in advance.

But in fixed-income markets, bond prices fluctuate randomly.
This means that discounting itself is risky.
Therefore, when pricing fixed-income products or long-dated derivatives,
it is not sufficient to model only the stock or underlying asset: one
must also model the dynamics of interest rates or bond prices.

\paragraph{Why this matters for derivative pricing.}

For short-dated equity options, assuming constant rates is often a
reasonable approximation.
But for products whose value depends strongly on the yield curve, such
as:
\begin{itemize}[leftmargin=1.4em]
    \item bonds,
    \item bond options,
    \item caplets and swaptions,
    \item long-dated structured products,
\end{itemize}
the dynamics of interest rates become essential.

In such cases, the discount factor is no longer deterministic, and the
pricing problem must explicitly incorporate the evolution of the term
structure.

\paragraph{Need for richer models.}

These observations motivate stochastic interest-rate models, in which
one models objects such as:
\begin{itemize}[leftmargin=1.4em]
    \item the short rate \(r_t\),
    \item bond prices \(P(t,T)\),
    \item yields,
    \item forward rates.
\end{itemize}

The goal is to build models that are flexible enough to describe the
observed yield curve and its random evolution through time.

\paragraph{Main lesson.}

The assumption of constant interest rates is useful as a first
approximation, but it is not realistic enough for many applications.
Once interest-rate risk matters, the whole term structure must be
treated as a genuine source of randomness, and this leads naturally to
stochastic models for fixed-income markets.
\subsection{Discount factors and zero-coupon bonds}

We begin with the most basic objects of fixed-income modeling.

\paragraph{Discount factors.}
If the short rate is denoted by \(r_t\), the value at time \(t\) of one
unit of cash received at time \(T\) is represented by the discount
factor
\[
D(t,T):=\exp\!\Big(-\int_t^T r_s\,\mathrm{d}s\Big).
\]

When the short rate is deterministic, this quantity is known at time
\(t\).
When interest rates are stochastic, the discount factor itself becomes
random, which is one of the main reasons why fixed-income pricing is
more subtle than in the constant-rate Black--Scholes setting.

\paragraph{Zero-coupon bonds.}
A zero-coupon bond maturing at time \(T\) is a contract that pays one
unit of currency at time \(T\) and nothing before.
Its price at time \(t\le T\) is denoted by
\[
P(t,T).
\]

Thus \(P(t,T)\) is the market price at time \(t\) of one unit of cash
delivered at time \(T\).

\paragraph{Basic properties.}
By definition,
\[
P(T,T)=1.
\]
Moreover, for fixed \(t\), the function
\[
T\longmapsto P(t,T)
\]
typically decreases with maturity: one unit of currency delivered far in
the future is worth less than one unit delivered soon.

\paragraph{Interpretation.}
Zero-coupon bonds are the elementary building blocks of fixed-income
markets.
Once the full family
\[
\{P(t,T):T\ge t\}
\]
is known, many other quantities of interest can be derived from it, such
as yields, forward rates, and the prices of coupon-bearing bonds.

\paragraph{Special case of constant rates.}
If the short rate is constant and equal to \(r\), then
\[
P(t,T)=e^{-r(T-t)}.
\]
Thus the Black--Scholes discount factor is nothing but the price of a
zero-coupon bond in the constant-rate case.

\subsection{Spot rates, zero rates, and simply compounded rates}

Once zero-coupon bond prices are known, one can define several
equivalent ways of describing the term structure of interest rates.

\paragraph{Spot rate / short rate.}
The \emph{spot rate}, or \emph{short rate}, at time \(t\) is the
instantaneous interest rate prevailing at that date.
It is usually denoted by
\[
r_t.
\]

In continuous-time models, the money market account satisfies
\[
\mathrm{d}B_t=r_tB_t\,\mathrm{d}t.
\]

\paragraph{Zero rate.}
For a given maturity \(T>t\), the continuously compounded zero rate
between \(t\) and \(T\) is defined by the relation
\[
P(t,T)=e^{-R(t,T)(T-t)}.
\]
Equivalently,
\[
\boxed{
R(t,T)=-\frac{1}{T-t}\log P(t,T).
}
\]

Thus \(R(t,T)\) is the constant continuously compounded rate that would
produce exactly the same discount factor over the interval \([t,T]\).

\paragraph{Simply compounded rate.}
Another common market convention is the simply compounded rate
\(L(t,T)\), defined by
\[
P(t,T)=\frac{1}{1+L(t,T)(T-t)}.
\]
Equivalently,
\[
\boxed{
L(t,T)=\frac{1}{T-t}\Big(\frac{1}{P(t,T)}-1\Big).
}
\]

This convention is often used for money-market instruments and short
maturities.

\paragraph{Relation between conventions.}
The continuously compounded zero rate and the simply compounded rate
carry the same information, since both are determined by the bond price
\(P(t,T)\).
They are just different parametrizations of the same discount factor.

\paragraph{Interpretation.}
The key point is that the bond price \(P(t,T)\) is the primitive object,
while rates such as \(R(t,T)\) or \(L(t,T)\) are derived quantities.
Different markets prefer different conventions, but they all encode the
same underlying term structure.

\subsection{Yield curve and term structure of interest rates}

The \emph{term structure of interest rates} describes how interest rates
depend on the maturity.

\paragraph{Yield curve.}
At a given time \(t\), the yield curve is the function
\[
T\longmapsto R(t,T),
\qquad T\ge t,
\]
where \(R(t,T)\) is the zero rate associated with maturity \(T\).

Equivalently, one may describe the term structure through the bond price
curve
\[
T\longmapsto P(t,T).
\]

\paragraph{Why the yield curve matters.}
The yield curve summarizes how the market prices the time value of money
across maturities.
It is one of the most important objects in fixed-income markets.

In practice, the yield curve is used to:
\begin{itemize}[leftmargin=1.4em]
    \item discount future cash flows,
    \item value bonds and interest-rate derivatives,
    \item infer market expectations about future rates,
    \item compare borrowing and lending conditions across maturities.
\end{itemize}

\paragraph{Typical shapes.}
Empirically, the yield curve is not flat.
Depending on market conditions, it may be:
\begin{itemize}[leftmargin=1.4em]
    \item upward sloping,
    \item downward sloping,
    \item or hump-shaped.
\end{itemize}

This immediately shows why a single constant rate cannot be sufficient
to describe fixed-income markets.

\paragraph{Term structure as a family of prices.}
Mathematically, the term structure may be viewed as the family of bond
prices
\[
\{P(t,T):T\ge t\}.
\]
At each time \(t\), this is an entire curve in the maturity variable
\(T\), not just a single number.

This is one of the main conceptual differences between fixed-income
modeling and Black--Scholes-type equity modeling:
the state of the market is no longer described by one asset price only,
but by a whole curve.

\paragraph{Dynamic point of view.}
As time evolves, the term structure itself moves randomly.
Thus fixed-income modeling is concerned not only with the shape of the
curve at one date, but also with its stochastic evolution through time.

This motivates the introduction of short-rate models, factor models, and
forward-rate models.

\subsection{Forward rates}

A third useful way to describe the term structure is through
\emph{forward rates}.

\paragraph{Definition.}
The instantaneous forward rate \(f(t,T)\) is defined by
\[
\boxed{
f(t,T):=-\partial_T \log P(t,T).
}
\]

Equivalently, bond prices can be reconstructed from forward rates via
\[
\boxed{
P(t,T)=\exp\!\Big(-\int_t^T f(t,u)\,\mathrm{d}u\Big).
}
\]

Thus the forward-rate curve
\[
T\longmapsto f(t,T)
\]
contains exactly the same information as the bond-price curve or the
yield curve.

\paragraph{Interpretation.}
Roughly speaking, \(f(t,T)\) represents the instantaneous rate agreed at
time \(t\) for borrowing or lending over an infinitesimal future period
around maturity \(T\).

It should therefore be viewed as a market-implied future rate.

\paragraph{Connection with zero rates.}
Since
\[
P(t,T)=e^{-R(t,T)(T-t)},
\]
one has the relation
\[
R(t,T)
=
\frac{1}{T-t}\int_t^T f(t,u)\,\mathrm{d}u.
\]
Thus the zero rate is the average of the forward rates over the interval
\([t,T]\).

\paragraph{Why forward rates are important.}
Forward rates play a central role in modern fixed-income theory because:
\begin{itemize}[leftmargin=1.4em]
    \item they provide a local description of the yield curve,
    \item they are often more convenient to model than bond prices
    directly,
    \item they are the basic state variables in the
    Heath--Jarrow--Morton framework.
\end{itemize}

\paragraph{Simply compounded forward rates.}
In practice, one also uses forward rates over a finite accrual period
\([T_1,T_2]\), defined by
\[
\boxed{
L(t;T_1,T_2)
=
\frac{1}{T_2-T_1}
\Big(
\frac{P(t,T_1)}{P(t,T_2)}-1
\Big).
}
\]

This quantity is especially important in money-market and LIBOR-style
models.

\paragraph{Main lesson.}
Bond prices, zero rates, and forward rates are three equivalent ways of
describing the same term structure.
The choice between them depends mostly on convenience and on the type of
model or derivative one wants to study.
\subsection{Fixed-income securities}

We now introduce the main securities traded in fixed-income markets.
These instruments are built from future cash flows occurring at
different dates, and their valuation relies on the term structure of
interest rates.

\subsubsection{Zero-coupon bonds}

The most basic fixed-income security is the zero-coupon bond.

\paragraph{Definition.}
A zero-coupon bond maturing at time \(T\) is a contract that pays one
unit of currency at time \(T\) and nothing before.
Its price at time \(t\le T\) is denoted by
\[
P(t,T).
\]

Thus, at maturity,
\[
P(T,T)=1.
\]

\paragraph{Interpretation.}
The price \(P(t,T)\) represents the present value at time \(t\) of one
unit of currency delivered at time \(T\).
In this sense, zero-coupon bonds are the elementary building blocks of
fixed-income markets.

\paragraph{Special case of deterministic rates.}
If the short rate is deterministic, then
\[
P(t,T)=\exp\!\Big(-\int_t^T r_s\,\mathrm{d}s\Big).
\]
In the constant-rate case, this reduces to
\[
P(t,T)=e^{-r(T-t)}.
\]

\paragraph{Why zero-coupon bonds are fundamental.}
Once the whole family
\[
\{P(t,T):T\ge t\}
\]
is known, many other fixed-income securities can be priced by
decomposing their cash flows into combinations of zero-coupon bonds.

\subsubsection{Coupon-bearing bonds}

A coupon-bearing bond is a bond that pays intermediate coupons before
maturity, in addition to a final redemption payment.

\paragraph{Cash-flow structure.}
Suppose a bond pays coupons \(c_1,\dots,c_n\) at dates
\[
T_1<\cdots<T_n,
\]
and repays a final principal \(F\) at time \(T_n\).
Then the bond price at time \(t\le T_1\) is
\[
B_t
=
\sum_{i=1}^n c_i\,P(t,T_i)+F\,P(t,T_n).
\]

Thus a coupon-bearing bond can be viewed simply as a portfolio of
zero-coupon bonds.

\paragraph{Interpretation.}
This decomposition is extremely important:
the whole fixed-income market can be built from elementary discount
factors \(P(t,T)\).
In particular, once zero-coupon prices are known, pricing coupon bonds
is straightforward.

\paragraph{Yield to maturity.}
In practice, coupon bonds are often quoted by their
\emph{yield to maturity}, i.e.\ the constant rate that discounts all
future coupons and principal back to the current market price.
However, from a modeling point of view, the bond price itself is more
fundamental than its yield.

\subsubsection{Bootstrapping intuition}

In practice, the full zero-coupon curve is not directly observed.
Instead, market quotes are available for a collection of liquid
fixed-income instruments, such as deposits, futures, swaps, or coupon
bonds.

\paragraph{Basic idea.}
\emph{Bootstrapping} is the procedure used to recover the zero-coupon
bond prices
\[
P(0,T_1),\ P(0,T_2),\ \dots
\]
from these market instruments.

The principle is recursive:
\begin{itemize}[leftmargin=1.4em]
    \item use the shortest maturity instruments to determine the first
    discount factors;
    \item then use longer maturity instruments, subtracting the value of
    already known early cash flows;
    \item continue step by step to reconstruct the whole curve.
\end{itemize}

\paragraph{Simple example.}
Suppose a coupon bond with maturity \(T_2\) pays coupon \(c\) at
\(T_1\), and coupon plus principal \(c+1\) at \(T_2\).
If its market price is \(B_0\), then
\[
B_0=c\,P(0,T_1)+(c+1)P(0,T_2).
\]
If \(P(0,T_1)\) is already known from a shorter maturity instrument,
then \(P(0,T_2)\) can be deduced.

\paragraph{Why bootstrapping matters.}
Bootstrapping is the practical link between raw market quotes and the
term structure used in pricing models.
Before one can model rates dynamically, one must first construct the
initial yield curve consistently from market data.

\subsection{Factor models for the term structure}

The term structure of interest rates is an entire curve
\[
T\longmapsto P(t,T)
\qquad\text{or}\qquad
T\longmapsto f(t,T),
\]
so in principle it is infinite-dimensional.

\paragraph{Motivation.}
To obtain tractable models, one often assumes that the whole yield curve
is driven by a finite number of underlying random factors.
This leads to \emph{factor models} for the term structure.

\paragraph{General idea.}
A factor model assumes that bond prices, yields, or forward rates can be
expressed in terms of a finite-dimensional state process
\[
X_t=(X_t^{(1)},\dots,X_t^{(n)}).
\]
For example,
\[
P(t,T)=F(t,T,X_t)
\]
for some deterministic function \(F\).

\paragraph{Interpretation.}
The factors \(X_t\) may represent:
\begin{itemize}[leftmargin=1.4em]
    \item the overall level of rates,
    \item the slope of the curve,
    \item the curvature of the curve,
    \item or more abstract latent variables.
\end{itemize}

\paragraph{Why factor models are useful.}
They reduce the complexity of the full term structure to a finite number
of state variables, which makes:
\begin{itemize}[leftmargin=1.4em]
    \item simulation easier,
    \item calibration more tractable,
    \item and derivative pricing more manageable.
\end{itemize}

\paragraph{Examples.}
Short-rate models are the simplest examples of factor models, where the
entire term structure is driven by the short rate alone or by a small
number of state variables.
More general frameworks, such as affine term-structure models, also fit
into this factor-model viewpoint.

\subsection{Short-rate models}

One of the most classical approaches to interest-rate modeling is to
model directly the short rate
\[
r_t.
\]

\paragraph{General principle.}
Once the short rate is specified as a stochastic process, one defines
bond prices by discounting future cash flows:
\[
P(t,T)
=
\E^\Qbb\!\left[
\exp\!\Big(-\int_t^T r_s\,\mathrm{d}s\Big)
\;\middle|\;
\mathcal{F}_t
\right].
\]

Thus the entire term structure is derived from the dynamics of the short
rate.

\paragraph{Examples.}
Typical short-rate models include:
\begin{itemize}[leftmargin=1.4em]
    \item the Vasicek model,
    \[
    \mathrm{d}r_t = a(b-r_t)\,\mathrm{d}t+\sigma\,\mathrm{d}W_t,
    \]
    \item the Cox--Ingersoll--Ross (CIR) model,
    \[
    \mathrm{d}r_t = a(b-r_t)\,\mathrm{d}t+\sigma\sqrt{r_t}\,\mathrm{d}W_t.
    \]
\end{itemize}

\paragraph{Interpretation.}
These models aim to capture key qualitative features of interest rates,
such as:
\begin{itemize}[leftmargin=1.4em]
    \item mean reversion,
    \item randomness in the short rate,
    \item and, in some cases, positivity.
\end{itemize}

\paragraph{Advantages.}
Short-rate models are conceptually simple and often analytically
tractable.
In many cases, one can derive explicit or semi-explicit formulas for
bond prices.

\paragraph{Limitations.}
However, they also have important drawbacks:
\begin{itemize}[leftmargin=1.4em]
    \item they model only the short rate directly,
    \item the whole yield curve is generated indirectly,
    \item and fitting the observed initial term structure may require
    additional adjustments.
\end{itemize}

\subsection{The Heath--Jarrow--Morton framework}

A more general approach is to model directly the evolution of the
forward-rate curve.

\paragraph{Basic idea.}
Instead of specifying the short rate \(r_t\), the
Heath--Jarrow--Morton (HJM) framework models the instantaneous forward
rates
\[
f(t,T)
\]
for all maturities \(T\ge t\).

Thus the primitive object is the family
\[
\{f(t,T):T\ge t\},
\]
rather than the short rate alone.

\paragraph{Why this is natural.}
Since the term structure is fundamentally a curve, it is natural to
model its evolution directly.
This makes HJM a very flexible and conceptually appealing framework.

\paragraph{Generic dynamics.}
One typically writes
\[
\mathrm{d}f(t,T)=\alpha(t,T)\,\mathrm{d}t+\sigma(t,T)\,\mathrm{d}W_t.
\]

The key point is that the drift \(\alpha(t,T)\) cannot be chosen
arbitrarily.
Under no-arbitrage, it is completely determined by the volatility
structure \(\sigma(t,T)\).

\paragraph{HJM drift restriction.}
This is the central feature of the HJM framework:
\begin{quote}
under the risk-neutral measure, once the volatility of the forward-rate
curve is specified, the drift is fixed by no-arbitrage.
\end{quote}

Thus the HJM framework provides a direct and general description of the
stochastic evolution of the entire term structure.

\paragraph{Advantages.}
The HJM approach is very flexible and naturally adapted to the
modeling of bond markets and interest-rate derivatives.

\paragraph{Drawbacks.}
Its main drawback is that it is infinite-dimensional in general, since
the whole forward-rate curve must be modeled.
To obtain tractable models, one often imposes additional structure,
for example finite-dimensional volatility specifications.

\subsection{Calibration issues in interest-rate models}

As in local volatility or stochastic volatility modeling, the practical
success of an interest-rate model depends crucially on calibration to
market data.

\paragraph{What must be fitted?}
An interest-rate model should typically reproduce:
\begin{itemize}[leftmargin=1.4em]
    \item the initial yield curve,
    \item the prices of liquid interest-rate derivatives,
    \item and, ideally, realistic future dynamics of the term structure.
\end{itemize}

\paragraph{Initial curve fit.}
A first requirement is exact fit to the observed initial term structure.
This is often achieved through bootstrapping and, in many models, by
adding a deterministic time-dependent shift.

\paragraph{Derivative calibration.}
A second requirement is consistency with market prices of liquid
fixed-income derivatives, such as:
\begin{itemize}[leftmargin=1.4em]
    \item caplets,
    \item swaptions,
    \item bond options.
\end{itemize}

These instruments provide information about the market's view of future
rate volatility.

\paragraph{Trade-off between realism and tractability.}
As always, there is a balance to strike:
\begin{itemize}[leftmargin=1.4em]
    \item simple models are easier to calibrate and use,
    \item richer models fit the market better, but may become
    numerically expensive or unstable.
\end{itemize}

\paragraph{Short-rate vs HJM calibration.}
Short-rate models are often easier to simulate and interpret, but may be
less flexible in fitting the whole curve and the derivative market.
HJM-type models are more flexible, but calibration is more demanding
because one must specify and estimate a whole volatility structure.

\paragraph{Main lesson.}
In interest-rate modeling, calibration is not only about finding
parameters.
It is also about choosing the right level of model complexity:
rich enough to fit the market, but simple enough to remain tractable and
stable in practice.
% ============================================================
\section{Practical comparison of the three approaches}
% ============================================================

We now step back and compare the three modeling directions introduced in
this lecture:
\begin{itemize}[leftmargin=1.4em]
    \item local volatility,
    \item stochastic volatility,
    \item stochastic interest-rate models.
\end{itemize}

Each of these frameworks is designed to improve on the Black--Scholes
model, but they do so in different ways and for different purposes.

\subsection{Local volatility vs stochastic volatility}

Local volatility and stochastic volatility models are both motivated by
the empirical failure of Black--Scholes to reproduce market smiles and
skews.

\paragraph{Local volatility.}
The local volatility model replaces the constant parameter \(\sigma\) by
a deterministic function
\[
\sigma_{\mathrm{loc}}(t,S_t).
\]
Its main strength is static calibration:
it can be constructed so as to reproduce the full cross-section of
European vanilla option prices observed at a given date.

Thus local volatility is very effective when the goal is to fit today's
market smile.

\paragraph{Stochastic volatility.}
By contrast, stochastic volatility models introduce an additional random
factor \(V_t\), so that volatility itself becomes random.
This gives a richer model, capable of producing more realistic
time-series behavior of volatility and more plausible future dynamics of
the implied volatility surface.

Thus stochastic volatility is more dynamic in nature:
it is not only a static fitting device, but also a model for how
volatility evolves through time.

\paragraph{Main difference.}
A useful summary is:
\begin{itemize}[leftmargin=1.4em]
    \item local volatility is mainly designed to fit the current
    volatility surface,
    \item stochastic volatility is mainly designed to model the future
    evolution of volatility.
\end{itemize}

\paragraph{Consequences for hedging and computation.}
Local volatility models remain one-factor and Markovian, so they are
closer in spirit to Black--Scholes and remain relatively tractable.

Stochastic volatility models are higher-dimensional, typically
incomplete, and more difficult to calibrate and hedge.
However, they often give a better description of market dynamics.

\subsection{When stochastic rates matter}

The need for stochastic interest rates depends strongly on the type of
product being priced.

\paragraph{Cases where constant rates may be sufficient.}
For short-dated equity options, assuming deterministic or constant rates
is often a reasonable approximation.
In such products, the main source of uncertainty comes from the
underlying asset, while the effect of rate dynamics is relatively small.

\paragraph{Cases where stochastic rates become important.}
By contrast, stochastic rates matter whenever:
\begin{itemize}[leftmargin=1.4em]
    \item the payoff depends directly on bond prices, yields, or forward
    rates,
    \item the maturity is long enough for rate uncertainty to become
    material,
    \item the product is written on fixed-income instruments,
    \item discounting itself is a significant source of risk.
\end{itemize}

In these situations, the whole term structure becomes a relevant
state variable, and modeling only the asset price is no longer enough.

\paragraph{Examples.}
Stochastic rates are essential for:
\begin{itemize}[leftmargin=1.4em]
    \item bonds and bond options,
    \item caplets, floorlets, and swaptions,
    \item long-dated structured products,
    \item hybrid products with both equity and interest-rate exposure.
\end{itemize}

\paragraph{Main lesson.}
The further one moves away from short-dated vanilla equity products and
toward fixed-income or long-dated derivatives, the more important the
stochastic modeling of interest rates becomes.

\subsection{What is calibrated from market data?}

A central practical question is: what object does each model attempt to
fit?

\paragraph{Local volatility calibration.}
In local volatility models, one calibrates the function
\[
\sigma_{\mathrm{loc}}(t,s)
\]
so as to match the observed prices of European vanilla options across
strikes and maturities.
Equivalently, one calibrates to the full implied volatility surface.

This is typically done through Dupire's formula or through numerical
optimization.

\paragraph{Stochastic volatility calibration.}
In stochastic volatility models, one calibrates a finite-dimensional set
of parameters, such as
\[
\kappa,\qquad \theta,\qquad \eta,\qquad \rho,\qquad V_0
\]
in the Heston model.

The calibration target is again the market implied volatility surface,
but the fit is obtained indirectly through these structural parameters.

Thus:
\begin{itemize}[leftmargin=1.4em]
    \item local volatility calibrates a whole function,
    \item stochastic volatility calibrates a finite set of parameters.
\end{itemize}

\paragraph{Interest-rate calibration.}
In interest-rate models, calibration usually has two layers:
\begin{itemize}[leftmargin=1.4em]
    \item fit the initial term structure of bond prices or yields,
    \item fit the prices of liquid rate derivatives, such as caplets or
    swaptions.
\end{itemize}

Hence the calibration targets are different from the equity setting:
instead of only fitting a vanilla option surface, one must also fit the
yield curve.

\paragraph{Main message.}
The calibration problem depends on the model class:
\begin{itemize}[leftmargin=1.4em]
    \item local volatility fits a volatility surface,
    \item stochastic volatility fits a smile through structural
    parameters,
    \item stochastic rates fit a yield curve and interest-rate
    derivative prices.
\end{itemize}

\subsection{Main modeling trade-offs}

All three approaches aim at improving realism beyond Black--Scholes, but
each comes with its own balance between flexibility, tractability, and
economic interpretation.

\paragraph{Local volatility.}
\begin{itemize}[leftmargin=1.4em]
    \item \textbf{Advantage:} excellent fit to today's vanilla market
    prices.
    \item \textbf{Drawback:} less realistic smile dynamics and no genuine
    volatility factor.
\end{itemize}

\paragraph{Stochastic volatility.}
\begin{itemize}[leftmargin=1.4em]
    \item \textbf{Advantage:} richer and more realistic dynamics,
    especially for volatility risk and smile evolution.
    \item \textbf{Drawback:} higher-dimensional models, incomplete
    markets, and more difficult calibration.
\end{itemize}

\paragraph{Stochastic interest rates.}
\begin{itemize}[leftmargin=1.4em]
    \item \textbf{Advantage:} realistic modeling of discounting and
    fixed-income markets.
    \item \textbf{Drawback:} the state variable is no longer just one
    asset price, but a full term structure or a factor representation of
    it.
\end{itemize}

\paragraph{General conclusion.}
There is no universally best model.
The appropriate choice depends on the product to be priced, the market
data available for calibration, and the balance one wants between:
\begin{itemize}[leftmargin=1.4em]
    \item exact fit to today's market,
    \item realistic future dynamics,
    \item computational tractability,
    \item and interpretability of the parameters.
\end{itemize}

% ============================================================
\section{Summary}
% ============================================================

\subsection{Main takeaways}

In this lecture, we introduced three major directions for going beyond
the Black--Scholes model.

The first one was the local volatility approach, where the constant
volatility parameter is replaced by a deterministic function of time and
spot.
This allows one to fit the whole implied volatility surface at a given
date and leads to Dupire's formula.

The second one was stochastic volatility, where volatility itself
becomes a random process.
This introduces an additional risk factor and leads to richer and more
realistic smile dynamics, at the cost of greater modeling and numerical
complexity.

The third one was stochastic interest-rate modeling, which becomes
necessary whenever the term structure of rates and discounting itself
play an important role.
This leads naturally to the study of bond prices, yields, forward rates,
factor models, and the HJM framework.

The main message is that all three extensions are motivated by the same
goal:
to improve realism beyond Black--Scholes while remaining sufficiently
tractable for pricing, calibration, and risk management.

\subsection{What students should remember}

After this lecture, students should remember the following key points.

\begin{itemize}[leftmargin=1.4em]
    \item \textbf{Local volatility.}
    A local volatility model has dynamics of the form
    \[
    \mathrm{d}S_t=rS_t\,\mathrm{d}t+\sigma_{\mathrm{loc}}(t,S_t)S_t\,\mathrm{d}W_t.
    \]
    It remains one-factor and Markovian.

    \item \textbf{Dupire's formula.}
    Local volatility can be recovered from European call prices across
    strikes and maturities through Dupire's formula.

    \item \textbf{Stochastic volatility.}
    A stochastic volatility model introduces an additional random factor
    \(V_t\), so the price process depends on both spot and volatility.

    \item \textbf{Heston-type idea.}
    In models such as Heston, the parameters controlling mean reversion,
    volatility of volatility, and correlation play a central role in the
    shape of the smile and skew.

    \item \textbf{Fixed-income basics.}
    The main objects of interest-rate modeling are zero-coupon bonds,
    yields, forward rates, and the term structure of interest rates.

    \item \textbf{Factor and HJM models.}
    Interest-rate models may be built either from a finite number of
    factors, often through short-rate models, or directly through the
    forward-rate curve in the HJM framework.

    \item \textbf{Model choice.}
    Different models answer different needs:
    local volatility is strong for static calibration,
    stochastic volatility for dynamic smile modeling,
    and stochastic rates for rate-sensitive products and fixed-income
    markets.
\end{itemize}

\end{document}